\section{Avaliação}
\label{sec:avaliacao}

\subsection{Método de medição}
\label{sec:metodo}

As medições de carga correm num ambiente próprio, descrito em código no ficheiro
\texttt{loadtest/compose.yaml}: o
central e o PostgreSQL~16 com \textbf{2 vCPU e 4~GB cada}, e o Prometheus a recolher as métricas do central a
cada 5~s (\dec{4}). A máquina é um AMD Ryzen~7 3700X com 16 \emph{threads} e 31~GB de memória. Cada medição
recria o ambiente do zero, com a base de dados vazia e o central acabado de arrancar, para nenhuma herdar
dados nem \emph{caches} da anterior. As medições não correm no CI, porque os \emph{runners} partilhados dão
números com demasiado ruído.

\paragraph{Agentes simulados.}
O k6 simula 1\,000 \emph{hosts}, um por utilizador virtual, cada um com a sua chave e a sua ligação gRPC. Cada
um comporta-se como o \texttt{monitor-remote}: envia estatísticas a cada 30~s e erros a cada 10~s, em lotes
de até 100 registos, tenta de novo com \emph{backoff} nas falhas transitórias e desiste nas definitivas. Os
erros seguem o perfil da Tabela~\ref{tab:perfil-logs}, atribuído de forma determinística pelo índice do
\emph{host}. Simular os agentes em vez de correr 1\,000 processos \texttt{monitor-remote} é o que torna a
medição possível numa só máquina; a simulação usa o mesmo contrato (\texttt{.proto} do \texttt{monitor-api})
e as mesmas regras de envio.

\paragraph{Três fontes, cada uma para o que mede bem.}
O k6 dá a visão do agente: latência de cada pedido, atraso acrescentado pelo central (definido na
secção~\ref{sec:requisitos}) e lotes que o agente teria posto de parte. O Prometheus dá a do servidor:
percentis da ingestão e do \emph{dashboard}, CPU, memória e ocupação do \emph{pool} de ligações. A base de dados
dá a integridade: no fim de cada medição compara-se o número de registos confirmados ao agente com o número
gravado, e procuram-se duplicados. As contagens vêm do k6 e da base de dados, e não dos contadores do
Prometheus, que perdem o que acontece antes da primeira recolha.

\paragraph{Medições.}
\texttt{nominal} (\cen{SC1}) corre 10~minutos com o perfil de logs; \texttt{nominal-x2} e
\texttt{nominal-x5} repetem-no com a taxa de logs multiplicada por 2 e por 5, porque o perfil é uma hipótese
(análise de sensibilidade). \texttt{noisy-neighbour} (\cen{SC5}) põe 1\,\% dos \emph{hosts} a 200 logs/s durante
2~minutos, a meio de uma medição nominal. \texttt{recovery-storm} (\cen{SC2}) dá a cada \emph{host} 1~h de
\emph{backlog} (cerca de 441\,000 registos no total) e põe-nos todos a enviar ao mesmo tempo, enquanto o
\emph{dashboard} é consultado a cada 2~s. \texttt{capacidade-1}, \texttt{-2} e \texttt{-3} (\cen{SC3}) põem a mesma
carga acima da capacidade com 1, 2 e 3 instâncias atrás do nginx. \texttt{dashboard-30d} (\cen{SC4}) gera
primeiro 31 dias de histórico e corre a carga nominal por cima, com as consultas do \emph{dashboard}, em duas
fases: sem e com a retenção (secção~\ref{sec:sc4}). Desde \dec{9}, a retenção está ligada por omissão, por isso
as medições finais correm com ela ativa, ao contrário das iterações 3 a 5; sem dados expirados, cada execução
não apaga nada e custa poucos milissegundos.

\paragraph{Arquitetura e níveis de teste.}
A arquitetura define as unidades e as interfaces que cada nível de teste exercita
(secção~\ref{sec:documentacao}). A Tabela~\ref{tab:niveis-teste} diz em que nível se verifica cada grupo de
cenários.

\begin{table}[htbp]
  \centering\small
  \caption{Cenários de qualidade por nível de teste}
  \label{tab:niveis-teste}
  \begin{tabularx}{\linewidth}{L{0.2\linewidth}L{0.3\linewidth}Y}
    \toprule
    Nível & O que exercita & Cenários \\
    \midrule
    Unitário & uma unidade da vista de módulos, com \emph{fakes} das interfaces (\dec{14}) & \cen{TE2}, \cen{TE3}, \cen{MO1}, \cen{MO2}, \cen{AV3}, \cen{AV4}, \cen{MN2} \\
    Integração & as interfaces entre elementos: os repositórios num PostgreSQL real, o contrato (\texttt{buf}), o SDK contra o esquema do remote & \cen{AV2}, \cen{SC3} (liderança), \cen{SE1}, \cen{SE3}, \cen{MO2} \\
    Aceitação de QA: carga & o sistema inteiro sob carga, com os limites de recursos (k6) & \cen{SC1}--\cen{SC5}, \cen{IO1} \\
    Aceitação de QA: ponta a ponta & remote, SDK, nginx com TLS e central reais em Compose & \cen{AV1}, \cen{MN1}, \cen{SE1} \\
    Inspeção & \emph{checklist}, regras estáticas, experiência controlada & \cen{SE2} (ASVS), \cen{MO3}, camadas (\texttt{depguard}) \\
    \bottomrule
  \end{tabularx}
\end{table}

\paragraph{Limitações.}
O gerador de carga corre na mesma máquina que o sistema medido. Os limites de CPU do central e da base de
dados impedem que o k6 lhes tire recursos, mas não isolam totalmente a memória e o disco partilhados. Os
números absolutos dependem desta máquina; o que se compara são as iterações entre si, medidas da mesma
forma.

\subsection{Escalabilidade: antes e depois}
\label{sec:aval-escala}

As tabelas desta secção comparam duas iterações, medidas da mesma forma (secção~\ref{sec:metodo}): a
iteração~3, com o contrato gRPC e as métricas mas antes de qualquer decisão de escalabilidade (a
\emph{baseline} de carga), e a iteração~7, a versão final, com todas as decisões e a retenção ligada por
omissão. O texto cita também a iteração~4, a primeira medição depois de \dec{5} (escritas em lote) e \dec{6}
(controlo de admissão), para separar o efeito destas duas decisões do das seguintes.

\gerado{carga-sc1.tex}

\paragraph{Regime normal (\cen{SC1}).}
Já na \emph{baseline}, \cen{SC1} cumpria-se com uma folga muito grande (Tabela~\ref{tab:carga-sc1}): p95 da
ingestão de 6~ms contra um limite de 300~ms, e o central a usar 4\,\% dos 2 vCPU. O débito medido, 122 registos
por segundo, coincide com a estimativa feita antes de qualquer medição (cerca de 90 logs e 33 estatísticas por
segundo, secção~\ref{sec:estimativas}): a medição confirmou a estimativa e a conclusão que se tirou dela, de que
a ingestão nominal não era o problema de escala. Na versão final, o p95 é de 5~ms e o CPU de 3\,\%.

\gerado{carga-sensibilidade.tex}

\paragraph{Sensibilidade ao perfil de logs.}
Como o perfil de logs é uma hipótese, cada medição nominal repetiu-se com o dobro e o quíntuplo da taxa de logs
(Tabela~\ref{tab:carga-sensibilidade}). Na \emph{baseline}, com cinco vezes mais logs, o p95 subia para 10~ms e o
CPU para 8\,\%; na versão final ficam em 5~ms e 4\,\%, praticamente iguais ao regime normal, porque o custo de
um lote deixou de crescer com o número de registos. \cen{SC1} cumpre-se em todos os casos: a conclusão não
depende de a hipótese estar certa.

\gerado{carga-sc5.tex}

\paragraph{Vizinho ruidoso (\cen{SC5}).}
Enquanto 1\,\% dos \emph{hosts} envia 200 logs por segundo, o atraso p95 dos restantes foi de 6~ms na
\emph{baseline} e de 5~ms depois (Tabela~\ref{tab:carga-sc5}). O risco que motivava a parte de justiça entre
\emph{hosts} de \dec{6} (lacuna S10) não se materializou a esta escala. É o \emph{generate and test} a funcionar: a
medição refutou uma hipótese da análise, e essa parte de \dec{6} ficou desenhada mas não implementada
(secção~\ref{sec:d6}). Os próprios \emph{hosts} em rajada também ganharam: o seu atraso p95 desceu de 43 para 11~ms.

\gerado{carga-sc2.tex}

\paragraph{Recuperação (\cen{SC2}).}
Era o único cenário que falhava, e falhava onde a análise previa (Tabela~\ref{tab:carga-sc2}). Na
\emph{baseline}, os 1\,000 \emph{hosts} entregavam as 441\,000 linhas do \emph{backlog} em 1,4~minutos, sem perdas
nem duplicados, mas o \emph{dashboard} ficava inutilizável, com p95 de pelo menos 10~s (o limite superior do
histograma). O CPU do central nunca passou de 30\,\%; o recurso esgotado era o \emph{pool} de ligações à base de
dados (25 de 25), com mais de 11\,000 esperas por uma ligação livre, e 9,5\,\% das tentativas de envio
acabavam em \emph{timeout}.

Depois de \dec{5} e \dec{6} (iteração~4), o \emph{dashboard} passou a ter p95 de 262~ms durante a recuperação, e
\cen{SC2} passou a cumprir-se; na versão final, com o \emph{bulkhead} de \dec{9}, é de 145~ms. A ingestão nunca ocupa mais de 15 ligações (o limite de \dec{6}), e já não há esperas nem
\emph{timeouts}. Quase um quarto das tentativas (24,8\,\% na versão final) é rejeitado por sobrecarga, mas são
rejeições intencionais: o agente volta a enviar depois do atraso indicado, e as 441\,000 linhas chegam todas, sem
duplicados. O resultado menos óbvio é que a recuperação ficou \emph{mais rápida}, 42~s em vez de 1,4~minutos,
apesar das rejeições. Com escritas em lote cada pedido ocupa uma ligação durante muito menos tempo, e limitar quantos
pedidos competem pela base de dados ao mesmo tempo elimina a contenção que, antes, fazia cada um demorar
segundos. O \emph{trade-off} previsto em \dec{6} (recuperação mais lenta em troca de estabilidade) não se
verificou: as duas decisões juntas melhoraram ambos os critérios.

\paragraph{Regressão (iterações~5 e~7).}
Depois de \dec{8}, as medições de \cen{SC1}, \cen{SC2}, \cen{SC5} e da sensibilidade repetiram-se com uma
instância. Todos os cenários continuaram a cumprir-se, com valores dentro da variação entre execuções (o p95 do
\emph{dashboard} durante a recuperação foi de 380~ms, contra 262~ms na iteração~4; o limite é 1~s). A medição
final (iteração~7) repetiu todos os cenários de carga sobre a versão entregue, com a retenção ligada, e nenhum
regrediu. A retenção corre também durante \cen{SC2}, mas sem dados expirados cada execução custa
milissegundos (p95 de 9,5~ms).

\subsection{Escala horizontal (\cen{SC3})}
\label{sec:aval-sc3}

\gerado{carga-sc3.tex}

\cen{SC3} mede o débito com 1, 2 e 3 instâncias atrás do nginx, com a mesma carga: os 1\,000 \emph{hosts} enviam
lotes de 100 erros sem pausa (Tabela~\ref{tab:carga-sc3}). A parte de coordenação cumpre-se: em todas as
medições, cada \emph{job} teve um único líder e correu numa só instância, e no ensaio da secção~\ref{sec:d8} a
troca de líder demorou 12~s, contra um limite de 30~s. A parte do débito \textbf{não se cumpre}: duas e três
instâncias confirmaram 1,01 e 1,04 vezes os registos de uma na medição final (1,01 e 1,03 na iteração~5),
contra os 1,8 e 2,5 vezes pedidos.

A tabela mostra porquê. Cada instância do central usa em média 11\,\% ou menos do seu CPU; o que sobe com as
instâncias é a latência de cada pedido (p95 de 100, 199 e 312~ms), quase na proporção da concorrência admitida
(15, 30 e 45 pedidos). Pela lei de Little, débito igual com o triplo da concorrência e o triplo da latência quer dizer que os
pedidos passam a esperar num recurso partilhado: o PostgreSQL. Uma amostragem dos processos ativos do
PostgreSQL durante a carga confirma-o: 56\,\% das amostras estavam a executar, 26\,\% à espera do registo de
transações (WAL: inserção, escrita e sincronização) e 11\,\% em contenção nas mesmas páginas de dados. O limite
é o caminho de escrita da base de dados, e não o central.

A premissa de \cen{SC3}, uma carga acima da capacidade de uma instância, deixou de existir com \dec{5}. As
escritas em lote tornaram o central tão leve que o estrangulamento passou para a base de dados. Uma instância
confirmou cerca de 22\,000 registos por segundo no minuto de pico, cerca de 180 vezes os 122 registos por segundo
do regime normal. Este fator é um limite superior: foi medido com lotes cheios, o caso mais eficiente, e os
\emph{hosts} reais enviam lotes pequenos. Mesmo assim, a folga cobre com margem os 10\,000 \emph{hosts} do
\emph{backlog} do projeto (cerca de 1\,200 registos por segundo).

Por isso, a equipa decidiu não escalar a base de dados. Réplicas de leitura não aumentam o débito de escrita;
\emph{sharding} seria uma complexidade que nenhum cenário pede; e desligar a sincronização do WAL poria em risco
\cen{AV1}, porque o agente apaga do \emph{buffer} o que o central confirma. Escalar a base de dados fica como
evolução E9 (secção~\ref{sec:evolucao}), com um gatilho medido: débito sustentado acima de metade da capacidade aqui medida. \dec{8} mantém
o seu valor pela disponibilidade: o central deixa de ser um ponto único de falha, e um \emph{job} sobrevive à
morte da instância que o corria. É o \emph{generate and test} outra vez: a medição mostrou que a premissa de um
cenário deixou de ser verdadeira, e o relatório regista-o em vez de ajustar o alvo.

Na iteração~5, o \emph{dashboard} também piorava sob esta carga extrema (p95 de 4,9~s com uma instância e de
9,4~s com três): o limite de admissão de \dec{6} é por instância, e três instâncias metem três vezes mais
pedidos em simultâneo na base de dados. Na versão final fica entre 84 e 174~ms: em cada instância, a ingestão
(15 ligações) e o \emph{dashboard} (5, o \emph{bulkhead} de \dec{9}) têm compartimentos separados dentro do
\emph{pool} de 25, e as consultas, agora de custo limitado, libertam a ligação depressa. A regra de configuração mantém-se: com N instâncias, \texttt{INGEST\_MAX\_INFLIGHT} deve ser o orçamento
total da ingestão a dividir por N (ADR-0007).

\subsection{Trinta dias de dados (\cen{SC4})}
\label{sec:sc4}

\gerado{carga-sc4.tex}

A Tabela~\ref{tab:carga-sc4} compara a medição de \cen{SC4} antes de \dec{9} com a medição final, cada uma em
duas fases de 10~minutos sobre 31 dias de histórico (89~M estatísticas, 239~M erros, 143 a 150~GB). Antes, a
primeira fase corria sem retenção e a segunda com a retenção a 30 dias a cada minuto; na versão final a
retenção está ligada por omissão, e as duas fases são duas corridas com ela. \textbf{Antes, \cen{SC4} falhava nas duas partes.} As sete consultas do
\emph{dashboard} tinham p95 entre 29 e 74~s, e a ingestão, mesmo sem retenção, confirmava só cerca de 11\,000
dos 73\,000 registos esperados em 10~minutos, com 58\,\% dos pedidos rejeitados por sobrecarga e um atraso
acrescentado pelo central de 636~s no p95. Nenhum registo se perdeu: os que não entraram ficaram nos agentes,
por entregar, como \dec{19} e o ADR-0020 preveem. A causa foi isolada medindo cada consulta sozinha
(secção~\ref{sec:d9}): quatro consultas sem limite temporal ocupavam todo o \emph{pool} de ligações.

\textbf{Depois, \cen{SC4} cumpre-se.} Na medição final, o p95 das sete consultas fica entre 15 e 210~ms (o pior
é o das métricas de um mês de um \emph{host}; na iteração~6, a primeira medição depois de \dec{9}, era de
287~ms), e a ingestão sobre 30 dias de dados fica ao nível da medição nominal com a base vazia: p95 de 5~ms,
nenhum pedido rejeitado, todos os registos confirmados, sem duplicados. O CPU do PostgreSQL desceu de 67\,\%
para 7\,\% em média. Isoladas, as quatro consultas passaram de 55~s, 17~s, 13~s e
5,3~s para 19~ms, 0,6~ms, 2~ms e 265~ms (estas últimas com a \emph{cache} fria).

\paragraph{A retenção e a decisão de não particionar.}
Antes, cada execução da retenção demorava até 27~s, perto do \emph{timeout}, não pelo seu custo, mas porque
esperava pelas ligações que o \emph{dashboard} ocupava. Depois, cada execução demora 0,2 a 0,3~s em média
(0,48~s no p95, nas três corridas com retenção) e o número de registos expirados desce cerca de 733 mil por
hora: a retenção remove 2,6 vezes o que entra, e o atraso acumulado esvazia-se. O p95 da ingestão não subiu com
a retenção (na iteração~6, 7,4~ms sem ela e 4,9~ms com ela; a fase com retenção corre depois e encontra a
\emph{cache} quente, e a diferença é da ordem do ruído). Nenhuma das três condições escritas no ADR-0017 antes da medição se
verifica, por isso \textbf{o particionamento não se fez}. Fica como o passo seguinte, com um gatilho sobre
métricas que já existem: a duração da retenção (p95 acima de 15~s), o p95 da ingestão nos minutos em que a
retenção corre, e o número de registos com mais de 30 dias. É o resultado mais claro do \emph{generate and
test} deste trabalho: o plano previa particionar, e a medição mostrou que o problema estava noutro sítio.

\subsection{Disponibilidade e integridade dos dados}
\label{sec:av1-av4}

\paragraph{\cen{AV2}: \emph{host} sem resposta, com 1.000 \emph{hosts}.}
Um teste de integração, num PostgreSQL real, regista 1.000 \emph{hosts} com \emph{endpoints} num servidor
local; metade nunca responde, e as \emph{probes} esgotam os 5~s de \emph{timeout}. O servidor regista quando cada
\emph{host} recebe cada \emph{probe}. Antes de \dec{7}, o teste corre um ciclo; depois, corre o \emph{job} durante
3~minutos, como em produção. A Tabela~\ref{tab:av2} compara os dois (dados em
\texttt{docs/metrics/iter-6/av2-antes.txt} e \texttt{av2-depois.txt}).

\begin{table}[htbp]
  \centering\small
  \caption{\cen{AV2}: 1.000 \emph{hosts}, metade sem resposta, antes e depois de \dec{7}}
  \label{tab:av2}
  \begin{tabular}{lrrl}
    \toprule
    Medida & Antes & Depois & Alvo \\
    \midrule
    Intervalo entre \emph{probes} do mesmo \emph{host} (máximo) & 350,7~s & 35,4~s & $\le$ 45~s \Sim \\
    Pior atraso até ao alerta, além do limiar & 6~min~6~s & 50,4~s & $\le$ 60~s \Sim \\
    Tempo de cada \emph{host} em ``desconhecido'' por ciclo & 4~min~21~s & 0 & 0 \Sim \\
    Iteração mais longa do \emph{job} & 350,7~s & 20,2~s & $<$ 60~s \Sim \\
    \emph{Probes} por segundo (máximo em simultâneo) & --- & 32,4 (128) & --- \\
    \bottomrule
  \end{tabular}
\end{table}

A mediana do intervalo é 30,3~s e o máximo fica abaixo dos 38~s calculados. A iteração mais longa é a primeira,
com os 1.000 \emph{hosts} devidos ao mesmo tempo; nas seguintes, o \emph{jitter} já os espalhou. O pior atraso
soma ao intervalo máximo uma avaliação de alertas (15~s), sem contar com a fase entre a falha e a próxima
avaliação, que só o reduz. \cen{AV2} cumpre-se.

\paragraph{\cen{AV4}: rajada com o central em baixo.}
Um teste de integração no remote usa um \emph{buffer} SQLite real com um limite de 1~MiB (o rígido a 4~MiB) e
simula 12 ciclos de rajada com o central em baixo: em cada ciclo, uma aplicação escreve 1.000 INFO, 50 WARN e
20 ERROR, e o \emph{worker} de \dec{19} corre uma vez. Cada ciclo escreve metade do limite de uma só vez, muito
mais do que uma rajada real escreve em 10~s, para forçar descartes em quase todos os ciclos.
A Tabela~\ref{tab:av4} resume o resultado (dados em \texttt{docs/metrics/iter-6/av4.txt}).

\begin{table}[htbp]
  \centering
  \small
  \caption{\cen{AV4}: \emph{buffer} de 1~MiB, 12 ciclos de rajada com o central em baixo}
  \label{tab:av4}
  \begin{tabular}{lrl}
    \toprule
    Medida & Valor & Alvo \\
    \midrule
    Logs escritos (INFO / WARN / ERROR) & 12.000 / 600 / 240 & --- \\
    Tamanho depois de cada ciclo (máximo) & 0,88~MiB & $\le$ 1~MiB \Sim \\
    ERROR perdidos & 0 de 240 & 0 \Sim \\
    WARN perdidos & 0 de 600 & só depois dos INFO \Sim \\
    INFO descartados, contados e persistidos & 11.500 & todos contados \Sim \\
    INFO que restam & os 500 mais recentes & mais antigos primeiro \Sim \\
    \bottomrule
  \end{tabular}
\end{table}

O tamanho antes de cada verificação chegou a 1,36~MiB: é o excesso entre duas verificações, aqui exagerado de
propósito (meio limite por ciclo). Com o limite por omissão de 512~MiB e uma rajada real, o excesso estimado é
inferior a 1\,\%. Os contadores sobrevivem a reabrir o ficheiro, como num reinício do agente. \cen{AV4} cumpre-se.

\paragraph{\cen{AV3}: retoma sem reinício.}
Um teste no remote, com o relógio injetado e um central falso em memória, recusa a chave com
\texttt{PERMISSION\_DENIED}. A entrega pausa e, durante 4~min~59~s, nenhuma das duas categorias faz pedidos. O
operador corrige a causa; aos 5~minutos, a \emph{probe} envia um só registo, que leva o estado do \emph{buffer}
com a pausa assinalada, o central aceita-o, e a entrega retoma no mesmo ciclo e entrega tudo, sem reiniciar o
agente. Outros testes fixam que uma \emph{probe} falhada marca a seguinte para 5~minutos depois, que só há uma
\emph{probe} de cada vez, e que o \texttt{SIGHUP} troca a chave e faz a \emph{probe} de imediato. No processo
real, dois \texttt{SIGHUP} seguidos (chave nova e depois um \texttt{.env} sem chave) deixam o agente vivo,
com a chave nova no primeiro caso e a anterior no segundo (dados em \texttt{docs/metrics/iter-6/av3.txt}).
\cen{AV3} cumpre-se.

\paragraph{\cen{AV1}: central em baixo, ponta a ponta.}
O teste ponta a ponta de \dec{15} mata o central (\texttt{docker compose kill}, sem paragem graciosa), com o
balanceador de pé, e a aplicação escreve 500 logs com o SDK durante a falha. O remote continua a recolher
estatísticas e a tentar entregar, recebendo \texttt{UNAVAILABLE}. Quando o central volta, o condutor espera
que o \emph{buffer} esvazie, para o remote para ter um estado estável e compara os dois lados registo a
registo, pelo identificador de origem: tudo o que o remote marcou como entregue tem de estar no central, uma só
vez, e nada no central pode faltar no \emph{buffer}. A comparação abrange o \emph{buffer} inteiro, incluindo os
500 logs do caminho feliz entregues antes da falha. A Tabela~\ref{tab:av1} resume as execuções (dados em
\texttt{docs/metrics/iter-6/e2e/}).

\begin{table}[htbp]
  \centering\small
  \caption{\cen{AV1}: central morto, 500 logs gerados durante a falha}
  \label{tab:av1}
  \begin{tabular}{lrrl}
    \toprule
    Medida & 60~s em baixo & 1~h em baixo & Alvo \\
    \midrule
    Duração real da falha & 63~s & 3.603~s & --- \\
    Pendentes quando o central volta (logs / estatísticas) & 500 / 10 & 500 / 599 & --- \\
    Logs entregues e no central & 1.000 / 1.000 & 1.000 / 1.000 & --- \\
    Estatísticas entregues e no central & 18 / 18 & 634 / 634 & --- \\
    Perdidos & 0 & 0 & 0 \Sim \\
    Duplicados & 0 & 0 & 0 \Sim \\
    No central sem origem no \emph{buffer} & 0 & 0 & 0 \Sim \\
    Pendentes ou em quarentena no fim & 0 & 0 & 0 \Sim \\
    Tempo até entregar tudo, depois de o central voltar & 36,7~s & 202,0~s & --- \\
    \bottomrule
  \end{tabular}
\end{table}

O tempo de recuperação é o do \emph{backoff} do remote: depois de 60~s de falhas seguidas, a espera entre
tentativas já ia em cerca de 40~s (32 a 48~s, com a variação aleatória); o limite do \emph{backoff} é 5~minutos, que é por isso o
atraso máximo depois de uma falha longa. Com 1~h em baixo, a entrega completa demorou 202~s: até 5~minutos de
espera pela tentativa seguinte, e depois 500 logs e 599 estatísticas em lotes. As duas execuções foram feitas com
o código ainda por \emph{commit}; o \emph{workflow} de ponta a ponta (\texttt{e2e.yml}) repete-as a partir do
\texttt{main}. Os duplicados são impossíveis no central por construção (recibos por
host e identificador de origem); o que o teste prova é que o remote não cria identificadores novos ao repetir
um envio e que não marca como entregue nada que o central não tenha gravado. \cen{AV1} cumpre-se.

\subsection{Manutenção, testabilidade e interoperabilidade: MO1--MO3, TE2, TE3, IO1 e métricas de código}
\label{sec:mo1}

\paragraph{\cen{MO1}: uma métrica de alerta nova.}
Com o mecanismo de \dec{13} no lugar, implementámos uma métrica real que o projeto precisava: o número de
registos em quarentena reportado por cada agente (\texttt{buffer\_quarantined}, \dec{17}). A Tabela~\ref{tab:mo1}
compara o que a alteração tocou com a linha de base contada antes de \dec{13}: 7 sítios em 5 ficheiros (três
Go, um modelo HTML e a migração), a que se somam os dois ficheiros Go que qualquer migração obriga a mudar
(dados em \texttt{docs/metrics/iter-6/mo1.md}).

\begin{table}[htbp]
  \centering\small
  \caption{\cen{MO1}: o que uma métrica de alerta nova exige, antes e depois de \dec{13}}
  \label{tab:mo1}
  \begin{tabular}{lccl}
    \toprule
    Medida & Antes & Depois & Alvo \\
    \midrule
    Ficheiros Go com código da métrica & 3 & 1 & --- \\
    Ficheiros Go que registam a migração & 2 & 2 & --- \\
    Ficheiros Go de produção, no total & 5 & 3 & $\le$ 3 \Sim \\
    Migrações & 1 & 1 & 1 \Sim \\
    Modelos HTML & 1 & 0 & --- \\
    Sítios a mudar fora da métrica e da migração & 6 & 0 & 0 \Sim \\
    Camadas tocadas & 3 & 1 & --- \\
    \bottomrule
  \end{tabular}
\end{table}

Dos três ficheiros Go, só um é da métrica: a estratégia, com 36 linhas. Os outros dois registam a migração e
mudam em qualquer migração, de métricas ou não. Um executor que descobrisse as migrações pelo nome do ficheiro
tirá-los-ia, mas a ordem explícita das migrações é uma garantia que preferimos manter. Os testes tocados foram
dois: o do registo e o de caracterização, que passou a fixar também o valor da métrica nova numa base de
dados real. O teste que compara o registo com o \texttt{CHECK} da base de dados confirmou que a migração
acompanha a métrica. A alteração foi feita com um assistente de programação, em pouco mais de um minuto de
relógio; esse tempo não representa o de uma pessoa sem assistência, mas a dimensão da alteração (um ficheiro
novo a seguir um modelo que já existe) deixa folga larga para o alvo de uma pessoa-dia. \cen{MO1} cumpre-se.

\paragraph{\cen{MO2}: uma definição por repositório.}
A Tabela~\ref{tab:mo2} conta, antes e depois de \dec{12}, quantas vezes a lista dos níveis de log está escrita em
Go, por repositório (dados em \texttt{docs/metrics/iter-6/te2-te3-mo2.md}). As cópias que ficam no SQL e nas
traduções do \emph{protobuf} são cobertas por seis testes de consistência. Para confirmar que apanham uma
divergência, retirámos um nível do \texttt{CHECK} do central e, noutra vez, da lista do SDK: em ambos os casos o
teste respetivo falhou, e a alteração foi revertida. No remote, as duas validações do endereço do central
passaram a uma. \cen{MO2} cumpre-se.

\begin{table}[htbp]
  \centering\small
  \caption{\cen{MO2}: lista dos níveis de log escrita em Go, por repositório}
  \label{tab:mo2}
  \begin{tabular}{lccl}
    \toprule
    Repositório & Antes & Depois & Verificada contra \\
    \midrule
    central & 5 & 1 & \texttt{CHECK} do SQL, \emph{protobuf}, SQL das métricas \\
    remote & 1 & 1 & \texttt{CHECK} do \emph{buffer}, \emph{protobuf} \\
    SDK & 1 (privada) & 1 (exportada) & \texttt{CHECK} do \emph{buffer} \\
    hello-world & 1 & 0 & usa a do SDK \\
    \bottomrule
  \end{tabular}
\end{table}

\paragraph{\cen{MO3}: uma mudança de apresentação.}
Medimo-lo com uma experiência controlada, numa cópia do central descartada no fim (dados e \emph{diff} em
\texttt{docs/metrics/iter-6/mo3.md}). Fizemos duas alterações de apresentação típicas, escolhidas antes de olhar
para o código que as serve. No \emph{dashboard}, os incidentes abertos passaram a ser o primeiro cartão, e o
cartão de erros passou a mostrar quantos dos erros de hoje são críticos. Na página de erros, o nível passou a
\emph{badge} com uma cor por severidade, e o ambiente ganhou uma coluna própria. As duas tocaram quatro
ficheiros, todos em \texttt{frontend/} (um modelo HTML por página, o JavaScript do \emph{dashboard} e a folha de
estilos), e nenhum ficheiro Go: \textbf{0 alterações em \texttt{services} e \texttt{repositories}}. Os testes
dos \emph{handlers}, que desenham os modelos reais, passaram sem alterações. \cen{MO3} cumpre-se. Não usámos
o histórico do git como medida porque os \emph{commits} são por iteração e juntam alterações de várias
naturezas. O limite da medida é o do próprio cenário: uma página que precise de um dado novo já não é só
apresentação, e toca nas camadas de baixo por desenho.

\paragraph{\cen{IO1}: um agente noutra linguagem.}
O gerador de carga do D10 é esse agente: está escrito em JavaScript, para o k6, e o único artefacto do sistema
que usa é o \texttt{.proto} do \texttt{monitor-api}, montado só de leitura (dados em
\texttt{docs/metrics/iter-6/io1.md}). Tudo o que precisa de saber sobre o comportamento está nos comentários do
contrato: a autenticação, a deduplicação por \texttt{source\_id}, o que fazer com cada código gRPC e com o
\texttt{RetryInfo}. A lógica de agente ocupa 101 linhas. Com as mesmas garantias do remote, confirmou 73\,229
registos na medição de \cen{SC4}, todos gravados, sem duplicados nem lotes perdidos, e fez o mesmo nos
441\,000 registos de \cen{SC2}. \cen{IO1} cumpre-se, com dois limites: a durabilidade local (o \emph{buffer}) é
responsabilidade de cada agente e não do contrato, e o agente foi escrito pela mesma equipa que escreveu o
contrato.

\paragraph{\cen{TE1}: integração contínua.}
Nos últimos 20 \emph{runs} com sucesso de cada repositório (lidos no GitHub a 21/11), o CI demorou em mediana
entre 89~s (\texttt{monitor-api}) e 141~s (\texttt{monitor-central}), e no máximo 401~s (\texttt{monitor-hello-world}),
abaixo dos 10~minutos (alvo). Uma falha bloqueia o \emph{merge}: o \texttt{main} de cada repositório está
protegido, também para os administradores, e aconteceu durante o trabalho. A Higiene do PR de \dec{5} e
\dec{6} falhou por falta de um \texttt{go mod tidy}, e o \emph{quality gate} do SonarQube recusou os PR do gRPC,
de \dec{4} e de \dec{5}--\dec{6} por cobertura do código novo abaixo de 80\,\%; nenhum entrou antes de ser
corrigido. \cen{TE1} cumpre-se.

\paragraph{\cen{TE2} e \cen{TE3}: \emph{handlers} e tempo.}
Com as interfaces de \dec{14}, a cobertura dos \emph{handlers} passou de 9,5\,\% para 95,6\,\% (alvo
$\ge$ 60\,\%). Os testes cobrem todos os \emph{handlers}, desenham as páginas reais e correm em cerca de 0,2~s,
sem base de dados. Nos testes dos quatro repositórios, os \texttt{time.Sleep} passaram de 1 para 0 (alvo 0). O
resumo diário é testado em seis casos com o relógio e o fuso controlados: antes, às e depois das 9~h, no fim do
mês, no dia em que a hora muda, e noutro fuso. \cen{TE2} e \cen{TE3} cumprem-se.

\gerado{metricas-codigo-antes-depois.tex}

\subsection{Monitorabilidade}
\label{sec:mn}

\paragraph{\cen{MN1}: quarentena visível no central.}
No remote, um teste põe um registo em quarentena com 150 estatísticas pendentes e verifica que os dois lotes do
ciclo levam a mesma fotografia, com a quarentena contada. No central, um teste de integração em PostgreSQL envia
um lote com três registos em quarentena e lê-os na API do \emph{dashboard}; verifica também que uma fotografia
mais antiga não substitui a mais recente e que uma inválida é ignorada sem perder as estatísticas. A quarentena
fica visível no lote seguinte.

O teste ponta a ponta de \dec{15} repete-o com os componentes reais: corrompe um log pendente no \emph{buffer}
(metadados que deixam de ser JSON), espera que o remote o ponha em quarentena e mede quanto tempo demora a
aparecer no central. Na primeira execução falhou, com 10,0~s para ciclos de 5~s: um ciclo de entrega que acorda
antes da recolha não tem estatísticas para levar a fotografia, e o central ficava dois ciclos atrás
(secção~\ref{sec:d15}). Depois da correção, três execuções mediram 5,85~s, 0,16~s e 5,82~s. O atraso depende da fase
entre a quarentena e a recolha seguinte, e o máximo é um ciclo de recolha, que no teste são 6~s: o intervalo
mais 1~s de amostragem do CPU. O limite do teste é 7~s, com 1~s para a observação pelo condutor; em produção,
o ciclo é de 31~s. \cen{MN1} cumpre-se, nos testes de cada lado e ponta a ponta.

\paragraph{\cen{MN2}: \emph{job} parado.}
Um teste corre um \emph{job} liderado com intervalo de 15~s que bloqueia, com o relógio controlado: aos 30~s o
\texttt{/health/ready} ainda responde 200, um segundo depois responde 503 e nomeia o \emph{job}, e o
\texttt{/health/live} continua a responder 200. Um \emph{job} liderado por outra instância nunca conta. Quando
a iteração termina, a instância volta a estar pronta. A deteção acontece em $\le$ 2 intervalos. \cen{MN2}
cumpre-se (dados em \texttt{docs/metrics/iter-6/mn1-mn2.txt}).

\subsection{Segurança}
\label{sec:aval-se}

\paragraph{\cen{SE1}: chave inválida ou revogada.}
O teste ponta a ponta arranca um segundo remote, com o seu próprio \emph{buffer}, duas vezes: com uma chave que
não existe e com a chave antiga de um host cuja chave foi regenerada (a mesma alteração que a interface faz:
um \emph{digest} novo na linha do host). Em cada caso, a aplicação escreve 10 logs. Nos dois casos, o remote
pausa por \texttt{UNAUTHENTICATED} em 0,7~s (\dec{16}); o central não tem nenhuma linha de logs, estatísticas
ou estado do \emph{buffer} de outro host que não o principal, e o host revogado nunca aparece como visto; no
\emph{buffer} do remote, os 10 logs e a estatística ficam pendentes, sem quarentena: corrigida a chave,
seriam entregues. \cen{SE1} cumpre-se.

\paragraph{\cen{SE2}: XSS e \emph{clickjacking}.}
Um teste verifica que todas as respostas, de página, de API, de erro e de rota inexistente, levam a CSP e os outros
cabeçalhos, e o teste das páginas falha se alguma tiver um \emph{script inline} ou de uma CDN. A prova principal
foi feita num navegador real, com a CSP em \emph{enforce}: as sete páginas da interface carregaram sem nenhum
relatório de violação; um \emph{script} injetado na página, como faria um XSS, não executou; e a página posta num
\texttt{<iframe>} foi bloqueada. Os dois bloqueios chegaram ao \emph{endpoint} de relatórios do central
(dados em \texttt{docs/metrics/iter-6/se2-se3.md}). \cen{SE2} cumpre-se.

\paragraph{\cen{SE3}: \emph{endpoint} de saúde para \emph{loopback} ou metadados.}
A guarda das \emph{probes} é testada com 14 endereços que têm de ser recusados (\emph{loopback}, metadados das
\emph{clouds}, \emph{link-local}, \emph{multicast}, endereços IPv4 escritos como IPv6) e 6 que têm de passar
(redes privadas e públicas). Num teste ponta a ponta, \emph{probes} reais para \texttt{127.0.0.1} e para
\texttt{localhost}, um nome que só depois de resolvido se sabe que é \emph{loopback}, não chegaram ao servidor
local, e os dois \emph{hosts} ficaram em baixo; sem a guarda, o mesmo servidor recebeu a \emph{probe}, o que mostra
que foi a guarda a travá-la. \cen{SE3} cumpre-se.

\paragraph{OWASP ASVS 5.0.0, nível 1.}
Dos 70 requisitos, antes de \dec{18} cumpriam-se 36, 25 não se aplicavam (não há \emph{upload} de ficheiros,
OAuth, \emph{tokens} autocontidos, XML nem WebSockets) e 9 falhavam; depois, cumprem-se 42 e falham 3, com o risco
aceite e um gatilho no ADR-0016 (avaliação completa em \texttt{docs/seguranca/asvs-n1.md}).

\subsection{Resumo dos cenários}
\label{sec:resumo-cenarios}

A Tabela~\ref{tab:resumo-cenarios} junta os 21 cenários do capítulo~\ref{sec:requisitos}, com a medida-alvo, o
valor medido e o veredicto. Os valores de carga são os da medição final (iteração~7); os restantes, das
medições indicadas em cada secção. Cumprem-se 20; \cen{SC3} cumpre-se na coordenação e não no débito, pela razão
da secção~\ref{sec:aval-sc3}.

\begin{table}[htbp]
  \centering\small
  \caption{Resumo dos cenários: alvo, valor medido e veredicto}
  \label{tab:resumo-cenarios}
  \begin{tabularx}{\linewidth}{lYYc}
    \toprule
    Cenário & Medida-alvo & Medido & Cumpre \\
    \midrule
    \cen{SC1} & p95 $<$ 300~ms; 0\,\% de erros; CPU $<$ 70\,\% & 5~ms; 0\,\%; 3\,\% & \Sim \\
    \cen{SC2} & \emph{backlog} em $\le$ 15~min; \emph{dashboard} p95 $<$ 1~s; 0 perdas e duplicados & 42~s; 145~ms; 0 e 0 & \Sim \\
    \cen{SC3} & débito $\ge$ 1,8× e $\ge$ 2,5×; 0 \emph{jobs} duplicados; troca de líder $\le$ 30~s & 1,01× e 1,04×; 0; 12~s & \Parc \\
    \cen{SC4} & \emph{dashboard} p95 $<$ 1~s com 30 dias; ingestão sem degradação com a retenção & 15 a 210~ms; ingestão p95 5~ms, 0 rejeições & \Sim \\
    \cen{SC5} & atraso p95 dos outros \emph{hosts} $<$ 5~s & 5~ms & \Sim \\
    \addlinespace
    \cen{AV1} & 1~h em baixo: 0 perdas, 0 duplicados & 3.603~s: 0 e 0 & \Sim \\
    \cen{AV2} & alerta em $\le$ \emph{timeout} + 60~s com 1\,000 \emph{hosts} & 50,4~s além do limiar & \Sim \\
    \cen{AV3} & retoma em $\le$ 5~min sem reiniciar & aos 5~min, sem reinício & \Sim \\
    \cen{AV4} & \emph{buffer} $\le$ limite; só INFO e WARN perdidos, contados & 0,88 de 1~MiB; 0 ERROR e 0 WARN perdidos; 11.500 INFO contados & \Sim \\
    \addlinespace
    \cen{MO1} & $\le$ 1 pessoa-dia; $\le$ 3 ficheiros Go e 1 migração & 3 ficheiros Go (1 da métrica) e 1 migração & \Sim \\
    \cen{MO2} & 1 sítio por repositório; divergência apanhada por teste & 1 por repositório; 2 divergências, 2 apanhadas & \Sim \\
    \cen{MO3} & 0 alterações em \texttt{services} e \texttt{repositories} & 0 (4 ficheiros, todos em \texttt{frontend/}) & \Sim \\
    \addlinespace
    \cen{TE1} & CI em $\le$ 10~min; uma falha bloqueia o \emph{merge} & mediana 89 a 141~s, máximo 401~s; bloqueou 3 PR & \Sim \\
    \cen{TE2} & cobertura dos \emph{handlers} $\ge$ 60\,\% & 95,6\,\% & \Sim \\
    \cen{TE3} & lógica temporal sem esperas & 0 \texttt{time.Sleep} nos testes & \Sim \\
    \addlinespace
    \cen{MN1} & quarentena visível em $\le$ 1 ciclo & no máximo 5,85~s, com ciclos de 6~s & \Sim \\
    \cen{MN2} & \texttt{/health/ready} falha em $\le$ 2 intervalos & 503 logo após 2 intervalos & \Sim \\
    \addlinespace
    \cen{SE1} & chave inválida ou revogada: recusada, nada persistido & recusadas as duas; 0 linhas no central & \Sim \\
    \cen{SE2} & XSS e \emph{clickjacking} bloqueados & \emph{script} injetado e \texttt{<iframe>} bloqueados & \Sim \\
    \cen{SE3} & \emph{probe} para \emph{loopback} ou metadados recusada & 14 de 14 endereços recusados & \Sim \\
    \addlinespace
    \cen{IO1} & agente noutra linguagem só com o contrato & agente k6 em JavaScript; 0 perdas e 0 duplicados & \Sim \\
    \bottomrule
  \end{tabularx}
  \par\smallskip\footnotesize \Sim{} cumpre-se; \Parc{} cumpre-se em parte.
\end{table}
